% 16 · The assignment
\section{THE ASSIGNMENT}
\label{sec:assignment}

\deckslides{slide 50, the hands-on brief.}

This chapter is the brief. Everything before it is the material you need to
execute it; everything you need to execute it is in this repository.

\subsection{What you are producing}

One object each (your assigned cluster) and one short report. The report makes
exactly three claims, each with evidence:

\begin{enumerate}
\item \textbf{A membership list} for your cluster, with the criterion stated
  precisely enough that another student can reproduce it: the cuts, the
  parameters, the random seed, the commit hash.
\item \textbf{An honest score} for that list, against both available referees
  (kinematic selection and Simbad memberships), reported as recall
  \emph{and} precision, and with the degenerate cases printed
  (\S\ref{sec:validation}).
\item \textbf{One physical parameter} derived from the list (an isochrone age
and distance, or a red-clump distance) with its uncertainty, and a sentence on
how the answer would change if your membership were 20\% larger or smaller.
\end{enumerate}

A report that claims a method is ``better'' without stating the population,
the task, the seed count and both halves of the metric pair has not met the
brief. This is the whole content of \S\ref{sec:validation}, applied to you.

\subsection{Setup}

The environment is reproducible from the repository root:

\begin{verbatim}
uv sync                     # python + dependencies
uv run cluster --help       # the CLI
uv run cluster download     # APOGEE DR19 + GALAH DR4 + Gaia DR3 subsets
\end{verbatim}

Field runs are the slow ones. For interactive work, set \texttt{CLUSTER\_FAST=1}
to cap the field sample; every number quoted in this workbook was produced with
\texttt{CLUSTER\_FAST=0}. A container with the same environment is provided
(\texttt{Dockerfile}), and the full walkthrough with worked examples is
\texttt{docs/student\_activities.md}.

\subsection{Your cluster}

\begin{table}[htbp]
\centering
\small
\tabcolsep4pt
\begin{tabular}{@{}llll@{}} \toprule cluster & kind & note & suggested track \\
\midrule M~3 & globular & metal-poor ($-1.6$) & A (spectral wins \\ M~5 &
globular & & A \\ M~13 & globular & & A \\ M~15 & globular & metal-poor
($-2.2$), southern & A or C \\ M~71 & globular & & A \\ M~107 & globular &
southern (DEC $-13^\circ$) & C \\ M~92 & globular & metal-poor ($-2.3$) & A \\
Pleiades & open & young ($0.1$ Gyr); chemistry fails here & A (expect a null)
\\ King~7 & open & sparse, faint & A \\ Berkeley~71 & open & sparse & A \\
IC~166 & open & sparse, contaminated & A \\ NGC~2158 & open & & A \\ NGC~1245 &
open & & A \\ King~5 & open & & A \\ NGC~7789 & open & giant-rich ($1.4$ Gyr) &
B) clump \\ NGC~1798 & open & & A \\ NGC~2420 & open & & A \\ NGC~6819 & open &
giant-rich; clump residual $+0.01$ & B \\ M~67 & open & the workhorse; most
data to beat & C (or A) \\ Berkeley~66 & open & sparse & A \\ NGC~188 & open &
old ($7.6$ Gyr) & B \\ NGC~6791 & open & old, metal-rich ($+0.4$) & B (or A) \\
Berkeley~17 & open & & A \\ NGC~2243 & open & the sweet spot: main sequence
\emph{and} clump & B or C \\ Collinder~261 & open & the sweet spot: 214 GALAH
members & C \\ \bottomrule
\end{tabular}
\caption{The 25 clusters and the track each one suggests. Track assignments
are advisory; the reasons are in the notes and in
\S\ref{sec:benchmark}--\S\ref{sec:physics}. Metadata for every cluster,
including the literature ages and distances you will be scored against, is in
\texttt{src/cluster/literature.py}.}
\label{tab:assignment}
\end{table}

Pair up on the open clusters; they are the hard part. Globulars tag cleanly
already, which makes them a faster and solo-friendly track, and it is worth
asking, while you work on one, why that is (\S\ref{sec:benchmark}: metallicity,
age and the absence of recent star formation all push in the same direction).

\subsection{The three tracks}

\textbf{Track A: spectral embeddings.} Compare the 16 abundances with the
256-dimensional spectral latent on your cluster:

\begin{verbatim}
uv run cluster run --spectral data/embeddings/attention_broad_merged.parquet
\end{verbatim}

Best on the metal-poor globulars, where element ratios lose discriminating
power and the spectrum retains it. Expect a null result on the Pleiades, and
report it as a result.

\textbf{Track B: isochrones and the red clump.} Recover age and distance:

\begin{verbatim}
uv run python scripts/red_clump.py
uv run python scripts/sweet_spot.py
\end{verbatim}

Best on the giant-rich open clusters (NGC~6819, NGC~7789). Compare against
\texttt{src/cluster/literature.py} and be explicit about the age--metallicity
degeneracy of \S 15.1.

\textbf{Track C: two surveys.} GALAH main sequence plus APOGEE giants:

\begin{verbatim}
uv run python scripts/build_galah_apogee.py
uv run python scripts/rerun_combined.py
\end{verbatim}

Southern clusters only (DEC $\lesssim +25^\circ$). This is the track with the
largest payoff and the largest number of ways to go quietly wrong: the parallax
relaxation is required, not optional, and the two surveys' abundance scales must
be reconciled before the combined fit means anything.

\subsection{Levers you may turn}

One axis at a time, and log the configuration each time: the sweep in
\textbf{Table~\ref{tab:levers}} was only interpretable because every run
recorded its settings.

\begin{table}[htbp]
\centering
\small
\tabcolsep5pt
\begin{tabular}{@{}lll@{}}
\toprule
lever & where & what it changes \\
\midrule
\texttt{TSNE.perplexity} & \texttt{src/cluster/config.py} & neighbourhood size of the embedding \\
\texttt{UMAP.n\_neighbors} & \texttt{src/cluster/config.py} & local graph scale \\
\texttt{HDBSCAN.min\_cluster\_size} & \texttt{src/cluster/config.py} & the size floor (removed by EVoC) \\
\texttt{SNR\_MIN} & environment & sample quality \\
\texttt{NORMALIZE\_ROWS} & environment & Euclidean $\leftrightarrow$ cosine equivalence \\
\texttt{USE\_ELEMENT\_WEIGHTS} & environment & element weighting, post-standardisation \\
\texttt{DWARF\_ONLY} & environment & giant/dwarf selection \\
\bottomrule
\end{tabular}
\caption{The levers, and the one rule that applies to all of them: change one
at a time, and report the effect with both halves of the metric pair. The
test-suite defaults are in \texttt{src/cluster/config.py}; the measured effects
on the DR17 sample are in \texttt{docs/experiment\_results.md}, and they do not
all survive on the current data.}
\label{tab:levers2}
\end{table}

\subsection{How you are assessed}

\begin{issues}[THE RUBRIC]
\textbf{Reproducibility (40\%).} A second person, on a clean checkout, gets
your numbers. Configuration, seed, commit hash and the exact command are in
your report. If you used a subsample, say how it was drawn.

\textbf{Empirical care (40\%).} Both referees. Both halves of recall and
precision. Mean $\pm$ s.d. over at least three seeds for any stochastic step.
The largest-group fraction printed. A null result reported as a null result: a
cluster that chemistry cannot separate is a finding, not a failure.

\textbf{Physics (20\%).} One parameter with an uncertainty, a comparison with
the literature value, and an honest statement of the dominant systematic, which
in most cases will be your membership list, not your photometry.
\end{issues}

\begin{issues}[EXERCISES]
\begin{enumerate}
\item Write the one-paragraph methods section of your report before you run
  anything: the population, the cuts, the method, the seed, the metrics, the
  referees. If you cannot write it, you are not ready to run.
\item Take a published tagging result in the field (a table of members, or a
claimed purity) and try to score it with this workbook's protocol. What
information is missing from the paper, and what would you have asked the
authors for?
\item Your cluster is listed with a suggested track. Write two sentences
  arguing for the other track, using only the cluster's properties (age,
  metallicity, declination, richness).
\end{enumerate}
\end{issues}
